\par\noindent\begin{minipage}{\linewidth}
\small
\setlength{\tabcolsep}{4pt}
\renewcommand{\arraystretch}{1.13}
\captionof{table}{Cambios entre períodos extremos y diferencias entre áreas.}\label{tab:S09}
\begin{tabular}{@{}>{\raggedright\arraybackslash}p{6.0cm}>{\raggedright\arraybackslash}p{2.4cm}>{\raggedright\arraybackslash}p{2.15cm}>{\raggedright\arraybackslash}p{2.15cm}>{\raggedright\arraybackslash}p{2.15cm}@{}}
\toprule
\textbf{Medida / diseño de candidatos} & \textbf{Cambio medio} & \textbf{Positivo} & \textbf{Negativo} & \textbf{No monótono} \\
\midrule
CKA, grupos propios & 0,0207 & 23 & 3 & 22 \\
CKA, 2.048 por grupo & 0,0206 & 24 & 2 & 22 \\
Rangos de distancias, propio & 0,0304 & 22 & 4 & 19 \\
Procrustes, propio & 0,0019 & 15 & 11 & 24 \\
k25, todos los candidatos & -0,0087 & 10 & 16 & 22 \\
k25, 2.048 / media & 0,0091 & 20 & 6 & 20 \\
k25, 2.048 / CLS & 0,0034 & 16 & 10 & 22 \\
k25, 2.048 / SEP & 0,0138 & 21 & 5 & 20 \\
\bottomrule
\end{tabular}
\end{minipage}\par
\par\smallskip{\small Final menos inicio: 2020--24 frente a 2000--04. Cada fila describe 26 áreas con igual peso; los cambios de vecinos son fracciones, por lo que se multiplican por 100 para obtener puntos porcentuales. Se usan cinco períodos para clasificar la monotonicidad. Los CSV incluyen controles k10/k25/k50, reglas, consultas idénticas y parejas. El número de candidatos puede cambiar el signo. Son contrastes descriptivos entre períodos, no convergencia histórica causal.\par}
